\subsubsection{把一組原子移到目標中心}

居中時，同一組中的所有原子使用同一個位移。設選擇中有 $N$ 個原子，分子用笛卡爾位置 $\mathbf r_i$，晶體用分數座標 $\mathbf f_i$，兩種情況分別為
\begin{equation}
 \begin{aligned}
 \Delta\mathbf r&=\mathbf r_{\mathrm{target}}-
        \frac1N\sum_{i=1}^{N}\mathbf r_i,
 &\mathbf r_i'&=\mathbf r_i+\Delta\mathbf r,\\
 \Delta\mathbf f&=\mathbf f_{\mathrm{target}}-
        \frac1N\sum_{i=1}^{N}\mathbf f_i,
 &\mathbf f_i'&=\mathbf f_i+\Delta\mathbf f.
 \end{aligned}
 \label{eq:center-selected-atoms}
\end{equation}
第一行用於分子，位置與位移的單位都是\AA；第二行用於晶體，各量均為無量綱分數座標。分子常用目標 $(0,0,0)$，晶體常用 $(0.5,0.5,0.5)$。每個原子在平均值中佔相同權重，因此居中不使用原子質量，也不改變所選原子之間的距離。

例如，分子中兩個原子的 $x$ 座標為1和3\,\AA，要把它們的平均位置移到原點，共同位移就是 $-2\,\text{\AA}$；處理後為 $-1$ 和1\,\AA，間距仍為2\,\AA。晶體中若兩個位點的第一分數分量為0.1和0.3，目標為0.5，則都加0.3，得到0.4和0.6。這裡只取一個分量說明計算，其餘兩個分量同樣處理。

週期邊界附近要先看清使用的是哪一個映像。0.99與0.01的普通平均為0.5，但兩個位點可能實際聚集在邊界兩側；若要以這段連續結構的中心為準，應先整理為相鄰的連續映像再居中。應用後坐標可能被歸回晶胞，因此還應從週期連線看整體，而不只比較某個座標的正負。

\subsubsection{繞指定錨點旋轉}

旋轉先從每個位置中減去錨點，再轉動這個相對向量，最後加回錨點。這樣，錨點本身不動，離錨點的距離也不變。沿用行向量表示，有
\begin{equation}
 \mathbf r'=\mathbf p+(\mathbf r-\mathbf p)
              R(\hat{\mathbf n},\theta)^{\mathsf T}.
 \label{eq:rotate-row-position}
\end{equation}
$\mathbf p$ 是以\AA 為單位的錨點，$\hat{\mathbf n}$ 是旋轉軸的單位方向，$\theta$ 是轉角；上標 $\mathsf T$ 表示轉置。對話方塊中填寫度數，正角遵循右手定則：右手拇指沿輸入軸的正向，四指彎曲的方向為正轉向。軸向量只決定方向，例如\literalinput{0 0 1}和\literalinput{0 0 2}表示同一根正向轉軸。

繞全域性 $z$ 軸旋轉時，常用的矩陣為
\[
 R_z(\theta)=\begin{bmatrix}
 \cos\theta&-\sin\theta&0\\
 \sin\theta&\cos\theta&0\\
 0&0&1
 \end{bmatrix}.
\]
取錨點 $\mathbf p=(1,1,0)\,\text{\AA}$，位置 $\mathbf r=(2,1,0)\,\text{\AA}$，繞 $z$ 軸正轉 $90^\circ$，得到 $\mathbf r'=(1,2,0)\,\text{\AA}$。這個小例子可用來辨認旋轉方向；旋轉後原子到錨點仍相距1\,\AA。晶體中原子的等價位置還會歸回晶胞，預覽時應連同邊界映像一起觀察。

\subsubsection{關於一個平面作映象}

映象需要一個平面上的點 $\mathbf p$ 和非零笛卡爾法向 $\mathbf n$。先把法向歸一化，再減去兩倍的有向投影：
\begin{equation}
 \begin{aligned}
 \hat{\mathbf n}&=\frac{\mathbf n}{\|\mathbf n\|},\\
 \mathbf r'&=\mathbf r-
 2\bigl[(\mathbf r-\mathbf p)\cdot\hat{\mathbf n}\bigr]
       \hat{\mathbf n}.
 \end{aligned}
 \label{eq:mirror-cartesian-plane}
\end{equation}
點積 $(\mathbf r-\mathbf p)\cdot\hat{\mathbf n}$ 是原子到平面的有向距離，單位為\AA。映象後它的符號反轉，大小相同；位於平面上的原子保持不動。法向填\literalinput{1 0 0}，平面點取 $(1,0,0)\,\text{\AA}$，就是 $x=1\,\text{\AA}$ 的平面。位置 $(1.5,2,0)\,\text{\AA}$ 映象後變為 $(0.5,2,0)\,\text{\AA}$，兩側到平面均為0.5\,\AA。

\begin{figure}[!htb]
\centering
\begin{tikzpicture}[x=1mm,y=1mm,font=\sffamily\small]
 \node[text=brand] at (24,47) {绕锚点旋转};
 \draw[teal,dashed] (10,15)--(33,15);
 \draw[teal,dashed] (10,15)--(10,38);
 \draw[->,>=stealth,brand,line width=.85pt] (33,15) arc[start angle=0,end angle=90,radius=23];
 \fill[ink] (10,15) circle (1.1);
 \fill[teal] (33,15) circle (1.4);
 \fill[brand] (10,38) circle (1.4);
 \node[anchor=north] at (10,13) {$\mathbf p$};
 \node[anchor=west,text=teal] at (35,15) {$\mathbf r$};
 \node[anchor=east,text=brand] at (8,38) {$\mathbf r'$};
 \node[text=brand] at (30,35) {$+90^\circ$};
 \node[text=muted,font=\sffamily\footnotesize] at (24,3) {到锚点的距离保持不变};
 \node[text=brand] at (86,47) {关于平面镜像};
 \fill[paleblue] (83,8) rectangle (87,40);
 \draw[teal,line width=.8pt] (85,8)--(85,40);
 \draw[dashed,muted] (68,29)--(102,29);
 \draw[->,>=stealth,brand,line width=.7pt] (85,12)--(105,12) node[above] {$\hat{\mathbf n}$};
 \fill[ink] (85,18) circle (1.1);
 \fill[teal] (102,29) circle (1.4);
 \fill[brand] (68,29) circle (1.4);
 \node[anchor=south,text=teal] at (102,31) {$\mathbf r$};
 \node[anchor=south,text=brand] at (68,31) {$\mathbf r'$};
 \node[anchor=east] at (83,18) {$\mathbf p$};
 \node[text=muted,font=\sffamily\footnotesize] at (86,3) {到镜面的距离两侧相等};
\end{tikzpicture}
\caption{锚点旋转与平面镜像（截面示意）。左图中的点绕 $\mathbf p$ 旋转；右图中的点跨过镜面，沿法向到镜面的距离大小不变。}
\label{fig:round10-anchor-mirror}
\end{figure}


居中、旋轉和映象均為幾何操作。處理分子時，先選擇需要整體移動的部分，再按下面的步驟操作；若只需調整觀察方向，旋轉相機即可。
